\section{Certification of the scalar inequalities}
\label{num:section}

We retain the definitions of $G$, $B$, and $S$ in
\eqref{nest:weight}, \eqref{nest:baseline}, and
\eqref{nest:gain-definition}. All terminating decimals in this appendix
denote exact rational numbers. In particular, the tabulated bounds below
are rational inequalities.

\begin{proposition}\label{num:scalar-positive}
The following strict inequalities hold:
\begin{align*}
 B(h)&>0 &&\text{for }h\in[.64,1.02]\cup[1.62,2.25],\\
 B(h)+S(h)&>0 &&\text{for }h\in[1.02,1.62].
\end{align*}
\end{proposition}

Finite knot bounds imply positivity on the whole curvature interval
through upper second derivative bounds that also control changes of formula.
The proof has three parts: a useful integral representation, those
derivative bounds, and finite rational certificates for interpolation.

\subsection{Layer cake and the baseline curvature bound}

Let $\rho$ be the inverse of the increasing function
$r\mapsto r/\sqrt{\pi(1-r)}$ on $[0,1)$. Explicitly, with $\sigma=\pi/4$,
\begin{equation}\label{num:rho}
 \rho(k)=\sqrt\pi\,k\sqrt{1+\sigma k^2}-\frac\pi2 k^2,
 \qquad k\ge0.
\end{equation}
Thus $t=\rho(1/h)$. On $[0,D]$ put $W(X)=-G'(X)$. From
\begin{equation}\label{num:weight-density}
 G(X)=\lambda\sqrt{1-X^2}+\frac{1-\lambda}{\sqrt{1-X^2}},
 \qquad
 W(X)=\frac{X(.96-.98X^2)}{(1-X^2)^{3/2}},
\end{equation}
we have $W\ge0$ and $\int_0^D W=1-\ell=.72$.

\begin{lemma}\label{num:layer-cake}
For $.64\le h\le2.25$,
\begin{equation}\label{num:layer-cake-formula}
 B(h)+1=\frac{\lambda h}{\pi}+\ell t+
 \int_0^D W(X)
 \min\left\{\pi(X/h)^2,\ X/h,\ \rho(X/h)\right\}\,dX.
\end{equation}
On an interval $[a,b]\subset[.64,2.25]$, the function
$B(h)-h^2/a^2$ is concave.
\end{lemma}

\begin{proof}
The reciprocal of the planar profile gives the increasing function
\[
 q_h(v):=\frac{hv}{p_*(v)}=
 \max\left\{h\sqrt{v/\pi},\ hv,
             \frac{hv}{\sqrt{\pi(1-v)}}\right\}.
\]
For $0\le X\le D<1$, its sublevel set in $[0,t]$ has length
\[
 \min\{\pi(X/h)^2,X/h,\rho(X/h)\}.
\]
Indeed these are the three separate sublevel restrictions; the additional
restriction $v\le t$ is redundant since
$\rho(X/h)<\rho(1/h)=t$. As a function of $k=X/h$, the selected
branches are $\pi k^2$ for $k\le a_0$, $k$ for
$a_0\le k\le b_0$, and $\rho(k)$ for $k\ge b_0$.
This is also obtained by inverting the three successive pieces of
$p_*$. The identity
\[
 G(q)=\ell+\int_0^D W(X)\mathbf 1_{\{q\le X\}}\,dX
 \qquad(q\ge0)
\]
and Tonelli's theorem now prove \eqref{num:layer-cake-formula}.

The function inverted to obtain $\rho$ has positive first and second
derivatives. Hence $\rho$ is increasing and concave, with $\rho(0)=0$,
and $k\rho'(k)\le\rho(k)\le1$. Consequently
\[
 \frac{d^2}{dh^2}\rho(X/h)
 =\frac{k^2\rho''(k)+2k\rho'(k)}{h^2}\le\frac2{h^2}.
\]
The same bound holds for $t=\rho(1/h)$. On the \emph{selected}
quadratic piece, $k\le a_0$, so
\[
 \frac{d^2}{dh^2}(\pi k^2)=\frac{6\pi k^2}{h^2}
 \le\frac6{\pi h^2}<\frac2{h^2}.
\]
On the selected linear piece the second derivative is
$2k/h^2<2/h^2$, since $k\le b_0<1$.

These bounds concern only the pieces that are actually selected.
For each fixed $X>0$, as $h$ increases they occur in the order
$\rho(X/h)$, $X/h$, $\pi(X/h)^2$. At the first change,
$\rho'(b_0)=2/(\pi+1)<1$, so the first derivative decreases.
At the second it changes from $-a_0/h$ to $-2a_0/h$,
again decreasing. There is thus no positive atom in the distributional
second derivative. The case $X=0$ is the constant zero function.
Each integrand, and also $t$, becomes concave on $[a,b]$ after
subtracting $h^2/a^2$. Their nonnegative weights in
\eqref{num:layer-cake-formula} sum to $\ell+.72=1$.
The remaining term is affine, proving the assertion.
\end{proof}

\subsection{The upper second derivative bound for the gain}

We say that a function has upper second derivative bound $K$ on an
interval if subtracting $Kh^2/2$ makes it concave there. This convention
includes changes of formula and does not require twice differentiability.

\begin{lemma}\label{num:s-second}
The function $S$ has upper second derivative bound $.4$ on
$[1.02,1.62]$.
\end{lemma}

\begin{proof}
Fix $y\ge0$, and write
\[
 F_y(h)=\ell y^2+C_0(h)f(w(h)),\qquad
 f(w)=\min\{d_m,w_+\}^3,
 \quad w=u-1+R-Qy,\quad R=2ht/\pi.
\]
Then $S=\inf_{y\ge0}F_y$. Differentiation of
$h^2t^2=\pi(1-t)$ gives
\begin{equation}\label{num:r-q-derivatives}
 \begin{gathered}
 R'=\frac{2t^2}{\pi(2-t)},\quad R''\le0,
 \qquad Q=\frac2{\sqrt\pi}(1+4h^2/\pi)^{1/4},\\
 \frac{Q'}Q=\frac{2h}{\pi+4h^2},\qquad
 \frac{Q''}Q=\frac{2\pi-4h^2}{(\pi+4h^2)^2}.
 \end{gathered}
\end{equation}
In particular, $0<Q'/Q\le1/(2h)$ and
$Q''/Q\ge-1/(4h^2)$.

Here are rational bounds sufficient for the argument. The inequalities
$3.14<\pi<22/7$ and the equation for $t$ imply
\[
 .64<t<.8,\qquad t(1.62)<.65.
\]
For example these follow by substituting the proposed endpoints into
$h^2t^2-\pi(1-t)$, which is strictly increasing in $t$.
Since $t$ decreases and $R$ increases, they yield
\[
 .19<\frac{1792}{9350}<R'
 <\frac{160}{471}<.34,
 \qquad
 R\le R(1.62)<\frac{1053}{1570}<.671.
\]
On a piece where $0<w<d_m$, it follows that
\[
 0\le Qy<P\le b_0-1+R<.36.
\]
The switch in $u=\min\{b_0,.96/h\}$ occurs at
$h_c=.96/b_0>1.4$. On its constant piece,
\[
 w'>.19-\frac{.36}{2(1.02)}>0,
 \qquad w''\le\frac{.36}{4(1.02)^2}<.8.
\]
On its reciprocal piece,
\[
 w'>-\frac{.96}{1.4^2}+.19-\frac{.36}{2(1.4)}>-.50,
 \qquad
 w''\le\frac{1.92}{1.4^3}+\frac{.36}{4(1.4)^2}<.8.
\]
On both pieces $w'\le R'<.34$. Also, on every smooth piece,
\[
 0<C_0\le.36,\qquad 0\le C_0'\le1.08,
 \qquad 0\le C_0''\le2.16;
\]
these follow directly from $C_0=\min\{.134h^3,.36\}$ and
$h\le1.62$.
Thus on an active cubic piece the product rule gives
\begin{align*}
 (C_0w^3)''
 &=C_0''w^3+6C_0'w^2w'+6C_0w(w')^2+3C_0w^2w''\\
 &\le2.16d_m^3+6(1.08)d_m^2(.34)
       +6(.36)d_m(.50)^2+3(.36)d_m^2(.8)\\
 &=.26866728<.4.
\end{align*}
On $w<0$ the added term is zero; on $w>d_m$ its second
derivative is at most $2.16d_m^3<.4$.

For completeness, all possible simultaneous changes of formula preserve
this bound. The only jumps in $w'$ and $C_0'$ are nonpositive,
coming respectively from the minima defining $u$ and $C_0$.
At a point with $0<w<d_m$ the jump in $F_y'$ is
\[
 w^3\Delta C_0'+3C_0w^2\Delta w'\le0.
\]
At $w=0$, both $f$ and its derivative vanish, so the derivatives
match even at a simultaneous switch. At $w=d_m$, the one-sided
derivatives of $f(w(h))$ are
\[
 3d_m^2\max\{w'_-,0\}\quad\hbox{from the left},\qquad
 3d_m^2\min\{w'_+,0\}\quad\hbox{from the right}.
\]
Their difference is nonpositive; a simultaneous switch of $C_0$
adds $d_m^3\Delta C_0'\le0$. These formulas include tangencies.
On a capped piece only the nonpositive jump of $C_0'$ remains.
For fixed $y$ these functions are algebraic on each branch, so the
threshold points are finite unless a threshold holds identically on a
piece; an identical threshold has the matching derivatives just described.
It follows that $F_y-.2h^2$ is concave on the full interval.
Finally,
\[
 S(h)-.2h^2=\inf_{y\ge0}\bigl(F_y(h)-.2h^2\bigr)
\]
is concave: the pointwise infimum of concave functions is concave,
as is seen either from its defining inequality or its hypograph.
The infimum is finite, since $F_y\ge0$ and $F_0\le C_0d_m^3$.
This also handles changes of the minimizing value of $y$.
\end{proof}

\subsection{Reduction of the knot values to arithmetic}

Only the factor $\rho(X/h)$ in \eqref{num:layer-cake-formula}
will require polynomial approximation. Set
$L=(D/h-b_0)_+$. If $L>0$, replace $\rho$ on $[b_0,D/h]$
by its cubic Taylor polynomial about $k_*=(b_0+D/h)/2$,
without changing either branch endpoint. Denote the resulting value
of $B$ by $\widetilde B$. If $L=0$, put $\widetilde B=B$.

\begin{lemma}\label{num:taylor-error}
For $.64\le h\le2.25$,
\begin{equation}\label{num:taylor-bound}
 B(h)\ge\widetilde B(h)-.008(D/h-b_0)_+^4.
\end{equation}
\end{lemma}

\begin{proof}
Differentiating \eqref{num:rho} gives
\[
 \rho'''(k)=\frac{3\sqrt\pi\,\sigma}{(1+\sigma k^2)^{5/2}},
 \qquad
 \rho''''(k)=-\frac{15\sqrt\pi\,\sigma^2 k}
                         {(1+\sigma k^2)^{7/2}}.
\]
The magnitude of the latter expression decreases when
$6\sigma k^2\ge1$, in particular for $k\ge b_0$.
Since $.68<b_0<.69$ and $3.14<\pi<22/7$,
\[
 |\rho''''(b_0)|
 <\frac{15\sqrt{22/7}(11/14)^2(.69)}
          {(1+(3.14/4)(.68)^2)^{7/2}}<4.
\]
The final inequality follows by squaring positive quantities and
cross-multiplying rational numbers. The midpoint Taylor remainder is
therefore at most $4(L/2)^4/24$ in absolute value. Multiplying
by the total weight $.72$ bounds its contribution to the integral
by $.0075L^4\le.008L^4$.
\end{proof}

Here are explicit moment formulas to evaluate $\widetilde B$.
For $0\le z\le D$, define
\[
 L_n(z)=\int_0^z\frac{X^n}{\sqrt{1-X^2}}\,dX,
 \qquad J_j(z)=\int_0^z X^jW(X)\,dX.
\]
Writing $\beta_z=\sqrt{1-z^2}$, integration by parts gives
\begin{equation}\label{num:moments-l}
 \begin{gathered}
 L_0(z)=\arcsin z,\qquad L_1(z)=1-\beta_z,\\
 L_n(z)=\frac{(n-1)L_{n-2}(z)-z^{n-1}\beta_z}{n}\qquad(n\ge2).
 \end{gathered}
\end{equation}
\begin{equation}\label{num:moments-j}
 J_j(z)=\lambda L_{j+1}(z)
 -(1-\lambda)\left(\frac{z^j}{\beta_z}
       -\mathbf1_{\{j=0\}}-jL_{j-1}(z)\right),
\end{equation}
where the term $jL_{j-1}$ is omitted for $j=0$.
Only $0\le j\le3$ and $0\le n\le4$ are needed. To check
\eqref{num:moments-j}, use \eqref{num:weight-density} and
differentiate $X^j/\sqrt{1-X^2}$; the indicator is its boundary
term at zero.

Put $x=ha_0$ and $y=\min\{D,hb_0\}$. Throughout the stated range,
$0<x\le y\le D$ and $x<D$. If the Taylor polynomial is
$\sum_{j=0}^3c_jk^j$, the evaluation formula is
\begin{equation}\label{num:evaluation-b}
 \widetilde B=\frac{\lambda h}{\pi}+\ell t-1
  +\frac\pi{h^2}J_2(x)+\frac{J_1(y)-J_1(x)}h
  +\sum_{j=0}^3\frac{c_j}{h^j}\bigl(J_j(D)-J_j(y)\bigr).
\end{equation}
The sum is omitted when $y=D$. Otherwise, put
$r_j=\rho^{(j)}(k_*)/j!$. Explicitly,
\begin{align*}
 r_0&=\sqrt\pi\,k_*\sqrt{1+\sigma k_*^2}-\pi k_*^2/2,\\
 r_1&=\frac{\sqrt\pi(1+2\sigma k_*^2)}{\sqrt{1+\sigma k_*^2}}
       -\pi k_*,\\
 r_2&=\frac12\left(
 \frac{\sqrt\pi\,\sigma k_*(3+2\sigma k_*^2)}
      {(1+\sigma k_*^2)^{3/2}}-\pi\right),
 &r_3&=\frac{\sqrt\pi\,\sigma}{2(1+\sigma k_*^2)^{5/2}},\\
 c_0&=r_0-k_*r_1+k_*^2r_2-k_*^3r_3,
 &c_1&=r_1-2k_*r_2+3k_*^2r_3,\\
 c_2&=r_2-3k_*r_3,&c_3&=r_3.
\end{align*}
The gain $S$ also has an evaluation formula using only arithmetic
and square roots.

\begin{lemma}[Evaluation of the capped gain]\label{nest:s-formula}
On $1.02\leq h\leq1.62$, $P,Q,C_0$ are positive.  With
\begin{equation}\label{nest:s-critical-point}
 q_*=\frac{2P}{1+\sqrt{1+6C_0Q^2P/\ell}},
\end{equation}
one has, including every cap and zero case,
\begin{equation}\label{nest:s-explicit}
 S(h)=\min\left\{C_0d_m^3,
       \ell\left(\frac{P-q_*}{Q}\right)^2+C_0q_*^3\right\}.
\end{equation}
\end{lemma}

\begin{proof}
Positivity of $Q,C_0$ is immediate.  To see that $P>0$ uniformly,
note that $t(h)$ decreases, so
$ht(h)=\sqrt{\pi(1-t(h))}$ increases.  At $h=1$ the equation
for $t$ gives $t(1)>3/4$, since $9/16-\pi/4<0$.
Using $h\geq1$, $u\geq16/27$, and $\pi<22/7$, we get
\[
 P>\frac{16}{27}-1+\frac{21}{44}
   =\frac{83}{1188}>0.
\]

For a fixed $h$, write $w=P-Qy$, so $w\leq P$. Define the
uncapped polynomial objective on $[0,P]$ by
\[
 F(w)=\frac{\ell(P-w)^2}{Q^2}+C_0w^3.
\]
It agrees with the original objective on the feasible uncapped interval
$[0,\min\{P,d_m\}]$. The polynomial is strictly convex on $[0,P]$,
with $F'(0)<0<F'(P)$, and its unique
critical point lies in $(0,P)$.  Solving
$3C_0Q^2w^2=2\ell(P-w)$ gives exactly $w=q_*$ in
\eqref{nest:s-critical-point}.

The branch $w\leq0$ has minimum $\ell P^2/Q^2=F(0)$,
which is greater than $F(q_*)$.  If $P\geq d_m$, the capped
branch $d_m\leq w\leq P$ has minimum $C_0d_m^3$, attained at
$w=P$.  When $q_*\leq d_m$ both candidates in
\eqref{nest:s-explicit} are consequently sufficient.  When
$q_*>d_m$, the minimum of $F$ on $[0,d_m]$ is $F(d_m)$,
which is at least $C_0d_m^3$; moreover
$F(q_*)\geq C_0q_*^3>C_0d_m^3$.  Thus the capped value still
wins in the displayed formula even though $q_*$ is infeasible.
Finally, if $P<d_m$, the cap is never reached, but its displayed
value is greater than $C_0P^3=F(P)\geq F(q_*)$ and hence cannot
alter the minimum.  The equalities $P=d_m$ and $q_*=d_m$
are included in these comparisons.  This proves the formula in
all cases.
\end{proof}

\subsection{An explicit outward arithmetic prescription}
\label{num:arithmetic}

We specify the arithmetic behind the table so that each entry can be
verified with integer operations. The input enclosure for $\pi$ is
\begin{equation}\label{num:pi-interval}
 \pi_-:=3.14159265358979<\pi<3.14159265358980=:\pi_+.
\end{equation}
Here is a rational verification of that enclosure. Set
\[
 A_N(z)=\sum_{j=0}^N\frac{(-1)^jz^{2j+1}}{2j+1},
 \qquad e_N(z)=\frac{z^{2N+3}}{2N+3}.
\]
For odd $N$ and $0<z<1$, the alternating series gives
$A_N(z)<\arctan z<A_N(z)+e_N(z)$. Machin's identity
$\pi=16\arctan(1/5)-4\arctan(1/239)$ follows by the tangent
addition formula, with the angle in $(0,\pi/2)$. Thus the following
entirely rational comparisons prove \eqref{num:pi-interval}:
\begin{align*}
 \pi_-&<16A_{31}(1/5)-4\bigl(A_9(1/239)+e_9(1/239)\bigr),\\
 16\bigl(A_{31}(1/5)+e_{31}(1/5)\bigr)-4A_9(1/239)&<\pi_+.
\end{align*}

The arcsine in \eqref{num:moments-l} can likewise be enclosed
without a transcendental evaluation. Define
\[
 m(z)=\frac{z}{\sqrt{2+2\sqrt{1-z^2}}},\qquad
 j_z=m(m(z)),\qquad
 T(j)=\sum_{n=0}^{10}\frac{\binom{2n}{n}}{4^n(2n+1)}j^{2n+1}.
\]
Twice applying the half-angle identity gives
$\arcsin z=4\arcsin j_z$. The function $m$ is increasing on
$[0,D]$, and
$j_D^2=(1-2/\sqrt7)/2<1/8<.36^2$, so $j_z<.36$.
The positive arcsine-series coefficients are at most one. Hence
\begin{equation}\label{num:arcsine-enclosure}
 4T(j_z)\le\arcsin z
 \le4T(j_z)+\frac{4j_z^{23}}{1-j_z^2}.
\end{equation}

For a precise common rounding rule, take $N=10^{40}$ and represent
an interval by integer endpoints $[a,b]_N=[a/N,b/N]$.
Enclose each rational input $r$ by
$[\lfloor Nr\rfloor,\lceil Nr\rceil]_N$. Addition and negation
are exact at this scale. For multiplication use
\[
 [a,b]_N[c,d]_N\subseteq
 \left[\left\lfloor\frac{\min\{ac,ad,bc,bd\}}N\right\rfloor,
       \left\lceil\frac{\max\{ac,ad,bc,bd\}}N\right\rceil\right]_N.
\]
For $0<a\le b$, reciprocal is enclosed by
$[\lfloor N^2/b\rfloor,\lceil N^2/a\rceil]_N$;
negative intervals are treated by negation. Division is multiplication
by this reciprocal. For $0\le a\le b$, square root is enclosed by
\[
 [\lfloor\sqrt{aN}\rfloor,\lceil\sqrt{bN}\rceil]_N,
\]
whose endpoints are found by integer-square comparisons.
The zeroth power is $[N,N]_N$.
An odd nonnegative integer power takes its extrema at the endpoints;
an even positive power has lower endpoint zero if the interval contains
zero, and otherwise the smaller endpoint power. Divide those integer
powers by $N^{n-1}$ and round outward. For minimum and maximum,
apply the same operation separately to the two lower and the two upper
endpoints. Each rule encloses every possible value of its operation,
also when its inputs are dependent. Induction therefore proves the
enclosure property for any expression formed by these rules.

Apply these rules to \eqref{num:evaluation-b} in the following order:
$\pi$, $a_0,b_0,D,t,x,y$; then $j_z$ and
\eqref{num:arcsine-enclosure} for $z=x,y,D$; then the recurrences
\eqref{num:moments-l}--\eqref{num:moments-j}; then $k_*,r_j,c_j$
if $hb_0<D$; and finally \eqref{num:evaluation-b}. Integer powers
are evaluated by the endpoint rule just given, and half-integer powers
as integer powers times a square root. Evaluate
\eqref{nest:s-explicit} from $t,u,P,Q,C_0,q_*$ in that order,
enclosing minima coordinatewise. No tabulated knot is a branch tie
$hb_0=D$, so that branch choice is determined by disjoint intervals.
In \eqref{num:arcsine-enclosure} retain the lower endpoint of
$4T$ and the upper endpoint of $4T+4j_z^{23}/(1-j_z^2)$.

\begin{table}[H]
\centering
\caption{Rational lower bounds: $10^5\widetilde B(h)>b_h$
and $10^5S(h)>s_h$. A dash means that $S$ is not used.}
\label{num:knot-table}
\begin{tabular}{rrr@{\qquad}rrr}
\toprule
$h$&$b_h$&$s_h$&$h$&$b_h$&$s_h$\\
\midrule
.64&1744&--&1.30&-196&271\\
.74&1200&--&1.34&-211&297\\
.84&735&--&1.38&-226&325\\
.92&432&--&1.42&-247&332\\
.98&247&--&1.46&-275&332\\
1.02&144&76&1.50&-260&318\\
1.06&57&98&1.54&-193&293\\
1.10&-15&123&1.58&-76&271\\
1.14&-74&151&1.62&86&251\\
1.18&-119&182&1.70&536&--\\
1.22&-153&215&1.90&2275&--\\
1.26&-178&247&2.25&6853&--\\
\bottomrule
\end{tabular}
\end{table}

Substitution by this prescription gives Table~\ref{num:knot-table}.
For additional checkable rounding slack, the resulting enclosures obey
\begin{equation}\label{num:table-slack}
 \begin{aligned}
 1.10&<10^5\widetilde B(h)-b_h<1.95
 &&\hbox{at all 24 knots},\\
 1.00&<10^5S(h)-s_h<1.66
 &&\hbox{at all 16 applicable knots}.
 \end{aligned}
\end{equation}
These inequalities include the arcsine-series tail and every arithmetic
rounding. They are deliberately coarser than the enclosures. For
example at the particularly small combined value $h=1.46$ they give,
more explicitly,
\[
 -273.170<10^5\widetilde B(1.46)<-273.169,
 \qquad 10^5S(1.46)=333.396.
\]
The latter equality follows from the cap term in
\eqref{nest:s-explicit}; here the Taylor error is zero.

\subsection{Certification between every pair of knots}

If $f$ has upper second derivative bound $K$ on $[a,b]$, concavity
of $f-Kh^2/2$ implies
\begin{equation}\label{num:chord}
 f(h)\ge\frac{b-h}{b-a}f(a)+\frac{h-a}{b-a}f(b)
             -\frac K2(h-a)(b-h)
 \ge\min\{f(a),f(b)\}-\frac{K(b-a)^2}{8}.
\end{equation}
By Lemmas~\ref{num:layer-cake} and~\ref{num:s-second}, use
$K=2/a^2$ for an exterior interval and $K=2.4$ for an interior
interval: indeed $2/(1.02)^2+.4<2.4$.

To make the final comparisons rational as well, set
\[
 D_+=.989743318611,\qquad
 b_-=1-1/\pi_-,\qquad
 E(h)=.008\max\{D_+/h-b_-,0\}^4.
\]
The integer comparison $49D_+^2>48$ shows $D<D_+$,
so $E(h)$ is an upper bound for the error in
\eqref{num:taylor-bound}. List the 24 knots of
Table~\ref{num:knot-table} in increasing order. For consecutive knots
$a,b$ set $\chi=1$ if $[a,b]\subset[1.02,1.62]$, and
$\chi=0$ otherwise. Then a rational uniform lower bound from
\eqref{num:chord} is
\begin{equation}\label{num:gap-formula}
 M_{a,b}=
 \min_{h\in\{a,b\}}
       \left\{\frac{b_h+\chi s_h}{10^5}-E(h)\right\}
       -\frac{K(b-a)^2}{8},
 \qquad K=\begin{cases}2.4,&\chi=1,\\2/a^2,&\chi=0.\end{cases}
\end{equation}
Here $s_h$ is omitted when $\chi=0$, including at the shared endpoints
$1.02$ and $1.62$. Every entry in Table~\ref{num:gap-table} is
verified directly from \eqref{num:gap-formula} by rational
cross-multiplication. Thus it requires no integral or transcendental
evaluation in addition to Table~\ref{num:knot-table}.

\begin{table}[H]
\centering
\caption{All 23 interpolation intervals. The last column of each pair
is an integer lower bound for $10^6M_{a,b}$.}
\label{num:gap-table}
\begin{tabular}{lr@{\qquad}lr}
\toprule
$[a,b]$&$10^6M_{a,b}\ge$&$[a,b]$&$10^6M_{a,b}\ge$\\
\midrule
$[.64,.74]$&4416&$[1.26,1.30]$&209\\
$[.74,.84]$&2298&$[1.30,1.34]$&269\\
$[.84,.92]$&1859&$[1.34,1.38]$&379\\
$[.92,.98]$&1313&$[1.38,1.42]$&369\\
$[.98,1.02]$&967&$[1.42,1.46]$&90\\
$[1.02,1.06]$&1037&$[1.46,1.50]$&90\\
$[1.06,1.10]$&581&$[1.50,1.54]$&100\\
$[1.10,1.14]$&280&$[1.54,1.58]$&520\\
$[1.14,1.18]$&145&$[1.58,1.62]$&1470\\
$[1.18,1.22]$&137&$[1.62,1.70]$&250\\
$[1.22,1.26]$&137&$[1.70,1.90]$&1899\\
&&$[1.90,2.25]$&14266\\
\bottomrule
\end{tabular}
\end{table}

\begin{proof}[Proof of Proposition~\ref{num:scalar-positive}]
Apply \eqref{num:chord} with the corrected knot lower bounds in
\eqref{num:gap-formula}. Table~\ref{num:gap-table} gives a strictly
positive lower bound on every interval, and those intervals cover
the stated ranges, including all endpoints. The smallest certified
uniform margin is $9/100000$, on $[1.42,1.46]$ and $[1.46,1.50]$.
In particular, the corrected lower bound at $1.46$ is exactly
$.00057$ and the interior interpolation loss is $.00048$.
\end{proof}
